\chapter[C'ye Giriş ve Kontrol Akışı · \textit{Temel}]{C'ye Giriş ve Kontrol Akışı}
\chaptermark{C'ye Giriş ve Kontrol Akışı}

Bölüm 20'de algoritmik fikirleri sözde kod düzeyinde tasarlamayı ele almıştık. Ancak bu fikirleri donanımın anlayacağı, sınır durumları hatasız işleyen ve zaman kısıtına uyan bir koda dönüştürmek sürecin tamamlayıcı adımıdır. C dili, donanım kaynaklarına doğrudan erişim sağlaması, bellek modelini açık biçimde sunması ve minimum çalışma zamanı (runtime) ek yükü getirmesi nedeniyle TÜBİTAK 1. Aşama sınavının ve yarışma programlamasının temel araçlarındandır.

Bu doğrultuda C programlama dilinin temel yapı taşlarını oluşturan veri tipleri, değişkenlerin bellekteki temsili, standart girdi/çıktı işlemleri ve program akışını belirleyen koşullu dallanma mekanizmaları aşağıda adım adım açıklanmıştır.

\section{Değişkenler, Tipler, printf/scanf}

Bir algoritmanın çalışabilmesi için veriyi bellekte saklaması, dönüştürmesi ve girdi/çıktı kanallarıyla etkileşime girmesi gerekir. C dilinde her verinin donanım seviyesinde karşılık gelen bir boyutu, yorumlanma biçimi (tipi) ve bellekte belirli bir adresi bulunur.

\subsection{İlk Program ve Bellek Sezgisi}

Temel bir programla başlayarak işletim sistemi ile kod arasındaki çalışma ilişkisini görelim:

\begin{lstlisting}[language=CTemel]
#include <stdio.h>

int main(void) {
    printf("Tubitak Bilgisayar Olimpiyati\n");
    return 0;
}
\end{lstlisting}

Bu program derlenip çalıştırıldığında işletim sistemi \texttt{main} fonksiyonunu çağırır. \texttt{\#include <stdio.h>} yönergesi, standart girdi/çıktı (standard input/output) fonksiyonlarının bildirimlerini içeren başlık dosyasını koda dahil eder. \texttt{return 0;} ifadesi ise işletim sistemine programın başarıyla sonlandığını bildiren bir çıkış kodudur (exit status). Değerlendirme sistemlerinde \texttt{return 0;} haricinde sıfır olmayan bir değerle çıkmak veya programın çökmesi, \textit{Runtime Error} (Çalışma Zamanı Hatası) olarak kaydedilir.

\subsection{Temel Veri Tipleri ve Temsil Sınırları}

Bilgisayar belleği bayt dizilerinden oluşur. Bir değişken tanımlandığında derleyici, bu değişkene tipi kadar baytlık ardışık bellek alanı ayırır.

\begin{definition}[Temel Tam Sayı Tipleri]
C dilinde tam sayı değerlerini saklamak için kullanılan temel tipler ve 64-bit mimarilerdeki yaygın boyutları şunlardır:
\begin{itemize}
    \item \texttt{char}: 1 bayt (8 bit). Genellikle karakterleri saklar; sayısal olarak $[-128, 127]$ veya $[0, 255]$ aralığında bir tam sayıdır.
    \item \texttt{int}: Genellikle 4 bayt (32 bit). $[-2^{31}, 2^{31}-1]$ aralığındaki tam sayıları temsil eder. Sayısal karşılığı $[-2\,147\,483\,648, 2\,147\,483\,647]$ aralığıdır ($\approx \pm 2 \cdot 10^9$).
    \item \texttt{long long}: En az 8 bayt (64 bit). $[-2^{63}, 2^{63}-1]$ aralığındaki tam sayıları saklar ($\approx \pm 9.22 \cdot 10^{18}$).
    \item \texttt{unsigned}: İşaretsiz belirteci, en soldaki işaret bitini de büyüklük hesabına katarak aralığı $[0, 2^k - 1]$ yapar (örneğin \texttt{unsigned int} için $[0, 2^{32}-1] \approx [0, 4.29 \cdot 10^9]$).
\end{itemize}
\end{definition}

\begin{definition}[Kayan Noktalı Tipler]
Gerçel sayıları IEEE 754 standardına göre temsil eden tiplerdir:
\begin{itemize}
    \item \texttt{float}: Tek duyarlıklı (single precision), 4 bayt. Yaklaşık 7 anlamlı basamak hassasiyetine sahiptir.
    \item \texttt{double}: Çift duyarlıklı (double precision), 8 bayt. Yaklaşık 15--17 anlamlı basamak hassasiyetine sahiptir. Hassasiyet kaybı riskini azaltmak adına gerçel sayılarda genellikle \texttt{double} tercih edilir.
\end{itemize}
\end{definition}

\begin{theorem}[Tam Sayı Taşması ve Tanımsız Davranış]
C standartlarına göre, işaretli tam sayılarda (\texttt{signed int}, \texttt{signed long long}) değerin temsil sınırlarını aşması \textbf{tanımsız davranıştır} (undefined behavior, kısaca UB). İşaretsiz tam sayılarda (\texttt{unsigned}) ise taşma tanımsız değildir; Bölüm 14'te ele aldığımız modüler aritmetik kurallarına uygun biçimde hesaplama doğal olarak $2^k$ modülünde yürür:
\[
a + b \pmod{2^k}
\]
burada $k$ değişkenin bit genişliğidir.
\end{theorem}

\begin{proof}[Açıklama ve İspat Taslağı]
Derleyiciler optimizasyon adımlarında işaretli tam sayı taşmasının hiç gerçekleşmeyeceğini varsayar. Örneğin işaretli bir $x$ için \texttt{x + 1 > x} ifadesi doğrudan \texttt{1} (doğru) olarak sadeleştirilebilir; oysa $x = 2^{31}-1$ iken taşma olursa değer $-2^{31}$ olur. Derleyici bu taşma ihtimalini hesaba katmak zorunda olmadığından, kod optimize edilirken beklenmeyen mantıksal hatalar doğabilir. İşaretsiz türlerde ise taşma davranışı standart tarafından $2^k$ modülüyle güvence altına alınmıştır.
\end{proof}

\subsection{Biçimlendirilmiş Girdi ve Çıktı: printf ve scanf}

\texttt{printf} ve \texttt{scanf}, C dilinde giriş ve çıkış işlemlerinin merkezinde yer alır. Her iki fonksiyon da biçim dizgileri (format strings) ve yer tutucular üzerinden çalışır.

\begin{definition}[Sık Kullanılan Biçim Yer Tutucuları]
Temel tiplerin \texttt{printf} ve \texttt{scanf} ile okunup yazdırılmasında kullanılan belirteçler:
\begin{itemize}
    \item \texttt{\textbackslash\{\}\%d} veya \texttt{\textbackslash\{\}\%i}: İşaretli 32-bit tam sayı (\texttt{int}).
    \item \texttt{\textbackslash\{\}\%u}: İşaretsiz 32-bit tam sayı (\texttt{unsigned int}).
    \item \texttt{\textbackslash\{\}\%lld}: İşaretli 64-bit tam sayı (\texttt{long long}).
    \item \texttt{\textbackslash\{\}\%llu}: İşaretsiz 64-bit tam sayı (\texttt{unsigned long long}).
    \item \texttt{\textbackslash\{\}\%f}: \texttt{printf} içinde \texttt{float} veya \texttt{double}, \texttt{scanf} içinde yalnızca \texttt{float}.
    \item \texttt{\textbackslash\{\}\%lf}: \texttt{scanf} içinde \texttt{double} okumak için zorunludur. \texttt{printf} içinde \texttt{double} için hem \texttt{\textbackslash\{\}\%f} hem de \texttt{\textbackslash\{\}\%lf} geçerlidir.
    \item \texttt{\textbackslash\{\}\%c}: Tek bir karakter (\texttt{char}).
    \item \texttt{\textbackslash\{\}\%s}: Boşluk karakteriyle sonlanan karakter dizisi (string).
\end{itemize}
\end{definition}

\texttt{scanf} fonksiyonuna argüman aktarırken temel kural şudur: Değişkenin doğrudan değeri değil, bellekteki adresi iletilmelidir. Bu işlem adres operatörü olan \texttt{\&} (ampersand) ile yapılır:
\begin{lstlisting}[language=CTemel]
int x;
scanf("%d", &x); // x'in bellek adresine girilen sayiyi yazar
\end{lstlisting}

\subsection{Çözümlü Örnekler}

\begin{example}
Klavyeden girilen iki adet $32$-bit tam sayıyı okuyup toplamlarını ekrana yazdıran, ancak toplam $32$-bit sınırını aştığında dahi doğru sonucu veren bir C programı yazınız.
\end{example}

\begin{proof}[Çözüm]
Girdiler $a$ ve $b$ olsun. Problemde $a, b \in [-2\cdot 10^9, 2\cdot 10^9]$ aralığında değerler alabilir. Bu durumda toplam uç noktalarda yaklaşık $\pm 4\cdot 10^9$ seviyesine ulaşır ve $32$-bit işaretli \texttt{int} türünün $[-2^{31}, 2^{31}-1]$ sınırını aşar.

Bu nedenle toplama işlemi yapılmadan önce değerlerden en az biri \texttt{long long} türüne dönüştürülmelidir. Aksi halde \texttt{int c = a + b;} ifadesinde toplama işlemi önce \texttt{int} genişliğinde yapılır, taşma gerçekleşir (UB) ve sonuç hatalı olur.

Program şu şekilde yazılabilir:
\begin{lstlisting}[language=CTemel]
#include <stdio.h>

int main(void) {
    int a, b;
    if (scanf("%d %d", &a, &b) != 2) {
        return 0;
    }
    long long toplam = (long long)a + b;
    printf("%lld\n", toplam);
    return 0;
}
\end{lstlisting}
Burada \texttt{(long long)a} ifadesi açık tip dönüşümü (type casting) gerçekleştirir. C'nin tür dönüşüm kurallarına göre bir terim \texttt{long long} olduğunda, diğer terim $b$ de otomatik olarak bu türe yükseltilir ve toplama işlemi 64-bit genişliğinde taşmasız tamamlanır.
\end{proof}

\begin{example}
Aşağıdaki C kod parçası çalıştırıldığında çıktısı ne olur? Nedenini Bölüm 15'te gördüğümüz ikiye tümleyen (two's complement) gösterimi üzerinden açıklayınız.

\begin{lstlisting}[language=CTemel]
unsigned int u = 0;
u = u - 1;
printf("%u\n", u);
\end{lstlisting}
\end{example}

\begin{proof}[Çözüm]
İşaretsiz 32-bit tam sayı türündeki bir değişken $0$ değerindeyken bellekteki 32 bitin tamamı sıfırdır:
\[
(00000000\ 00000000\ 00000000\ 00000000)_2
\]
Bu değerden $1$ çıkarıldığında işaretsiz taşma kuralı devreye girer. Aritmetik $2^{32}$ modülünde çalıştığından:
\[
0 - 1 \equiv -1 \equiv 2^{32} - 1 \pmod{2^{32}}
\]
elde edilir. Ortaya çıkan bit dizilimi $32$ adet $1$ içerir:
\[
(11111111\ 11111111\ 11111111\ 11111111)_2 = 2^{32} - 1 = 4\,294\,967\,295
\]
Böylece ekrana \texttt{4294967295} yazdırılır.
\end{proof}

\begin{example}
Aşağıdaki kod parçasının ekrana ne yazdıracağını belirleyiniz:
\begin{lstlisting}[language=CTemel]
int a = 7, b = 2;
double d1 = a / b;
double d2 = (double)a / b;
double d3 = (double)(a / b);
printf("%.2f %.2f %.2f\n", d1, d2, d3);
\end{lstlisting}
\end{example}

\begin{proof}[Çözüm]
İfadeler adım adım incelendiğinde:
\begin{enumerate}
    \item \texttt{a / b}: İki değişken de \texttt{int} türündedir. C'de tam sayı bölmesinde kesirli kısım doğrudan atılır; dolayısıyla $7 / 2 = 3$ elde edilir. Bu $3$ değeri \texttt{double} türündeki \texttt{d1} değişkenine aktarılırken $3.0$ olur.
    \item \texttt{(double)a / b}: Tür dönüşümünün işlem önceliği yüksektir. \texttt{(double)a} ifadesi $7.0$ gerçel sayısını üretir. İşlem $7.0 / 2$ haline gelir; C'nin örtük tür yükseltme kuralı gereği $2$ de $2.0$ değerine çevrilir ve kayan noktalı bölme ile $3.5$ bulunur. \texttt{d2} değişkeni $3.5$ değerini alır.
    \item \texttt{(double)(a / b)}: Parantez içindeki \texttt{a / b} bölmesi önce yapılır ve tam sayı bölmesiyle $3$ sonucunu verir. Ardından bu değer \texttt{(double)} ile dönüştürülerek \texttt{d3} değişkenine $3.0$ olarak atanır.
\end{enumerate}
\texttt{\textbackslash\{\}\%.2f} belirteci virgülden sonra $2$ basamak yazdırdığı için çıktı:
\[
3.00\ 3.50\ 3.00
\]
olur.
\end{proof}

\begin{example}
\texttt{scanf} ile ardışık bir tam sayı ve bir karakter okunurken ortaya çıkan boşluk/satır sonu karakteri sorununu ve çözümünü gösteriniz.
\end{example}

\begin{proof}[Çözüm]
Şu kod parçasını ele alalım:
\begin{lstlisting}[language=CTemel]
int sayi;
char harf;
scanf("%d", &sayi);
scanf("%c", &harf);
\end{lstlisting}
Kullanıcı konsola \texttt{42} yazıp enter tuşuna bastığında, girdi akışına (stdin) sırasıyla \texttt{'4'}, \texttt{'2'} ve satır sonu belirten \texttt{'\textbackslash\{\}n'} karakteri eklenir. 
\begin{itemize}
    \item İlk \texttt{scanf("\textbackslash\{\}\%d", \&sayi)} çağrısı sayısal karakterleri okuyup \texttt{sayi = 42} yapar, ancak \texttt{'\textbackslash\{\}n'} karakterini akışta bırakır.
    \item İkinci \texttt{scanf("\textbackslash\{\}\%c", \&harf)} çağrısı ise \texttt{\textbackslash\{\}\%c} belirtecinin boşluk veya satır sonu karakterlerini otomatik atlamaması nedeniyle doğrudan akışta bekleyen \texttt{'\textbackslash\{\}n'} karakterini okur ve değişkene yazar. Program yeni bir karakter beklemeden akışı tamamlar.
\end{itemize}

Bu durumu önlemek için biçim dizesinde \texttt{\textbackslash\{\}\%c} önüne bir boşluk eklenir:
\begin{lstlisting}[language=CTemel]
scanf(" %c", &harf);
\end{lstlisting}
Baştaki bu boşluk karakteri, \texttt{scanf}'in girdi akışındaki tüm ardışık boşluk, sekme (\texttt{\textbackslash\{\}t}) ve yeni satır (\texttt{\textbackslash\{\}n}) karakterlerini tüketmesini sağlar.
\end{proof}

\subsection{En Sık Yapılan Hatalar}
\begin{itemize}
    \item \textbf{\texttt{scanf} içinde \texttt{\&} işaretini unutmak:} \texttt{scanf("\textbackslash\{\}\%d", x);} yazıldığında derleyici uyarı verebilse de kod derlenebilir. Program çalışma anında $x$'in mevcut değerini bir bellek adresi olarak kullanmaya kalkışır ve bellek erişim ihlali (\textit{Segmentation Fault}) ile sonlanır.
    \item \textbf{\texttt{double} okurken \texttt{\textbackslash\{\}\%f} kullanmak:} \texttt{printf} fonksiyonunda \texttt{double} için \texttt{\textbackslash\{\}\%f} yeterlidir; ancak \texttt{scanf} ile \texttt{double} okunurken mutlaka \texttt{\textbackslash\{\}\%lf} (long float) kullanılmalıdır. \texttt{\textbackslash\{\}\%f} kullanıldığında belleğe yalnızca 4 bayt yazılacağından veri bozulur.
    \item \textbf{İşaretli tam sayı taşmasını gözden kaçırmak:} Örneğin $10^5 \times 10^5 = 10^{10}$ işlemi ara adımda $32$-bit sınırını aşar. \texttt{long long res = a * b;} yazıldığında $a$ ve $b$ türü \texttt{int} ise çarpım önce $32$-bit olarak hesaplanır (ve taşar), ardından taşmış değer \texttt{res} içine atanır. Güvenli yaklaşım \texttt{(long long)a * b} biçiminde dönüşüm yapmaktır.
\end{itemize}

\newpage

\section{if-else ve switch}

Algoritmalar yalnızca ardışık adımlardan ibaret değildir; verinin durumuna göre farklı hesaplama yollarını seçmek gerekir. C dilinde bu karar mekanizmaları \texttt{if-else} blokları ve \texttt{switch} çoklu dallanma yapılarıyla kurulur.

\subsection{Mantıksal Doğruluk Sezgisi ve Karşılaştırma Operatörleri}

Bölüm 1'de mantıksal önermeleri ve doğruluk değerlerini görmüştük. C dilinin ilk standartlarında yerleşik bir boole (\texttt{bool}) tipi yer almıyordu; doğruluk mantığı doğrudan tam sayılar üzerinden tanımlanmıştır:
\begin{itemize}
    \item $0$ değeri \textbf{YANLIŞ} (false) kabul edilir.
    \item Sıfır dışındaki \textbf{bütün tam sayılar} (negatif sayılar dahil: $-1, 5, 100$) \textbf{DOĞRU} (true) kabul edilir.
\end{itemize}
Daha sonraki standartlarda \texttt{<stdbool.h>} başlığıyla \texttt{bool}, \texttt{true} ($1$) ve \texttt{false} ($0$) tanımlanmış olsa da arka plandaki sayısal kural değişmemiştir.

\begin{definition}[İlişkisel ve Mantıksal Operatörler]
C dilindeki temel karşılaştırma araçları:
\begin{itemize}
    \item İlişkisel: \texttt{==} (eşit mi), \texttt{!=} (farklı mı), \texttt{<}, \texttt{>}, \texttt{<=}, \texttt{>=}.
    \item Mantıksal: \texttt{\&\&} (mantıksal VE), \texttt{||} (mantıksal VEYA), \texttt{!} (mantıksal DEĞİL).
\end{itemize}
Bu operatörler, koşul sağlandığında tamsayı olarak \texttt{1}, sağlanmadığında \texttt{0} değeri üretir.
\end{definition}

\begin{theorem}[Kısa Devre Değerlendirmesi / Short-Circuit Evaluation]
Mantıksal ifadeler soldan sağa doğru incelenir. İfadenin genel sonucu kesinleştiği anda kalan kısım \textbf{çalıştırılmaz}:
\begin{enumerate}
    \item \texttt{A \&\& B} ifadesinde $A$ yanlış ($0$) ise, $B$'nin sonucuna bakılmaksızın ifadenin tamamı yanlıştır; dolayısıyla $B$ hiç hesaplanmaz.
    \item \texttt{A || B} ifadesinde $A$ doğru (sıfır dışı) ise, $B$'nin sonucuna bakılmaksızın ifade doğrudur; dolayısıyla $B$ hiç hesaplanmaz.
\end{enumerate}
\end{theorem}

\begin{proof}[Açıklama ve İspat Taslağı]
Kısa devre davranışı, derleyici tarafından örtük koşullu atlama komutlarına dönüştürülür:
{\small\[
\text{\texttt{if (A \textbackslash\{\}\&\textbackslash\{\}\& B)}} \iff \text{\texttt{if (A) \textbackslash\{\}\{ if (B) \textbackslash\{\}\{ ... \textbackslash\{\}\} \textbackslash\{\}\}}}
\]}
Bu mekanizma tanımsız durumlara ve program çökmelerine karşı önemli bir güvencedir. Örneğin \texttt{b != 0 \&\& a / b > 2} kontrolünde $b=0$ olduğunda sağ taraftaki \texttt{a / b} bölmesi işletilmez; böylece sıfıra bölme hatası önlenir.
\end{proof}

\subsection{if-else Yapısı ve Asılı Else (Dangling Else) Problemi}

Temel koşullu dallanma kalıbı şu şekildedir:
\begin{lstlisting}[language=CTemel]
if (kosul1) {
    // kosul1 dogruysa calisir
} else if (kosul2) {
    // kosul1 yanlis ve kosul2 dogruysa calisir
} else {
    // hicbiri saglanmazsa calisir
}
\end{lstlisting}

Süslü parantez kullanılmadığında derleyici yalnızca ilk komutu blok kapsamına alır. Bu durum klasik \textit{Asılı Else} sorununu doğurabilir:
\begin{lstlisting}[language=CTemel]
if (x > 0)
    if (y > 0)
        printf("1");
else
    printf("2");
\end{lstlisting}
Buradaki girintileme yanıltıcıdır. C dilinin sözdizimi kuralına göre: \textbf{Bir \texttt{else}, kendisinden önce gelen ve henüz eşleşmemiş olan en yakın \texttt{if} bloğuna bağlanır.} Dolayısıyla kodun gerçek işleyişi şöyledir:
\begin{lstlisting}[language=CTemel]
if (x > 0) {
    if (y > 0) {
        printf("1");
    } else {
        printf("2");
    }
}
\end{lstlisting}
$x \le 0$ durumunda dış koşul sağlanmadığından ekrana herhangi bir çıktı verilmez.

\subsection{switch-case Yapısı ve Düşme (Fall-Through) Özelliği}

Bir tamsayı değişkenin çok sayıda sabit değere göre dallanması gerektiğinde \texttt{switch} yapısı tercih edilir:

\begin{lstlisting}[language=CTemel]
switch (ifade) {
    case sabit1:
        // ifade == sabit1 ise buradan itibaren calisir
        break;
    case sabit2:
        // ifade == sabit2 ise buradan itibaren calisir
        break;
    default:
        // hicbir case tutmazsa calisir
        break;
}
\end{lstlisting}

\begin{definition}[Düşme (Fall-Through)]
Bir \texttt{case} bloğunun sonuna \texttt{break;} konulmadığında, program sonraki \texttt{case} etiketinin koşulunu denetlemeden alt satırları yürütmeye devam eder. Bu duruma \textit{düşme (fall-through)} adı verilir. Genellikle unutulan bir \texttt{break} nedeniyle hata kaynağı olsa da kimi zaman bilinçli olarak kullanılır.
\end{definition}

\subsection{Çözümlü Örnekler}

\begin{example}
Verilen bir yılın artık yıl (leap year) olup olmadığını belirleyen mantıksal ifadeyi ve C kodunu oluşturunuz. (Kural: Bir yıl 4 ile tam bölünüyorsa artık yıldır; ancak 100 ile tam bölünen yıllardan yalnızca 400 ile tam bölünenler artık yıldır).
\end{example}

\begin{proof}[Çözüm]
Girdi yılı $Y$ olsun. Şartları inceleyelim:
\begin{itemize}
    \item Durum 1: Yıl $400$'e bölünüyorsa koşulsuz artık yıldır ($Y \pmod{400} == 0$).
    \item Durum 2: Yıl $4$'e bölünüp aynı zamanda $100$'e bölünmüyorsa artık yıldır ($Y \pmod 4 == 0 \text{ ve } Y \pmod{100} \ne 0$).
\end{itemize}
Bu iki durum mantıksal VEYA (\texttt{||}) ile birleştirilebilir:
\[
(Y \pmod{400} == 0) \lor (Y \pmod 4 == 0 \land Y \pmod{100} \ne 0)
\]
Kısa devre kuralından yararlanmak için, yılların büyük bölümü 4'e bölünmediğinden \texttt{yil \% 4 == 0} kontrolünü en başa almak kontrol sürecini hızlandırır:

\begin{lstlisting}[language=CTemel]
#include <stdio.h>

int main(void) {
    int yil;
    if (scanf("%d", &yil) != 1) return 0;

    if ((yil % 4 == 0 && yil % 100 != 0) || (yil % 400 == 0)) {
        printf("%d artik yildir.\n", yil);
    } else {
        printf("%d artik yil degildir.\n", yil);
    }
    return 0;
}
\end{lstlisting}
Operatör önceliğinde \texttt{\&\&} operatörü \texttt{||} operatöründen önce gelse de açık parantez kullanımı kodun doğruluğunu ve okunabilirliğini korur.
\end{proof}

\begin{example}
Aşağıdaki C programı çalıştırıldığında ekrana ne basar?
\begin{lstlisting}[language=CTemel]
#include <stdio.h>

int main(void) {
    int x = 2;
    switch (x) {
        case 1:
            printf("A");
        case 2:
            printf("B");
        case 3:
            printf("C");
            break;
        case 4:
            printf("D");
            break;
        default:
            printf("E");
    }
    printf("\n");
    return 0;
}
\end{lstlisting}
\end{example}

\begin{proof}[Çözüm]
Program akışını adım adım takip edelim:
\begin{enumerate}
    \item \texttt{x} değişkeninin değeri $2$'dir.
    \item \texttt{switch} ifadesi doğrudan eşleşen \texttt{case 2:} etiketine dallanır.
    \item \texttt{printf("B");} çalışır ve ekrana \texttt{B} yazdırılır.
    \item \texttt{case 2:} bloğunun sonunda \texttt{break;} bulunmadığından akış bir sonraki satıra düşer (fall-through).
    \item \texttt{case 3:} etiketi denetlenmeksizin altındaki kod işletilir: \texttt{printf("C");} çalışır ve ekrana \texttt{C} eklenir.
    \item \texttt{case 3:} bloğunun sonundaki \texttt{break;} ifadesi \texttt{switch} gövdesini sonlandırır.
\end{enumerate}
Böylece ekrana basılan metin:
\[
\text{BC}
\]
olur.
\end{proof}

\begin{example}
Aşağıdaki kod parçasının çıktısını belirleyiniz:
\begin{lstlisting}[language=CTemel]
int a = 0, b = 5;
if (a++ && ++b) {
    printf("Dogru: %d %d\n", a, b);
} else {
    printf("Yanlis: %d %d\n", a, b);
}
\end{lstlisting}
\end{example}

\begin{proof}[Çözüm]
Bu soru kısa devre değerlendirmesi ile artırma operatörlerinin etkileşimini sınar:
\begin{enumerate}
    \item \texttt{a++ \&\& ++b} ifadesinde önce sol taraf değerlendirilir: \texttt{a++}.
    \item \texttt{a++} ifadesinin o anki değeri $a$'nın başlangıçtaki değeri olan $0$'dır. İfade değerlendirildikten hemen sonra yan etki (side-effect) olarak $a$ bir artırılarak $1$ olur.
    \item Mantıksal VE (\texttt{\&\&}) işleminin sol tarafı $0$ (yanlış) değerini aldığından kısa devre kuralı devreye girer.
    \item Sağ taraftaki \texttt{++b} ifadesi hiç çalıştırılmaz; dolayısıyla $b$ değişkeni $5$ olarak kalır.
    \item Koşul sağlanmadığı için \texttt{else} bloğuna geçilir.
\end{enumerate}
Sonuç olarak ekrana basılan çıktı:
\[
\text{Yanlis: 1 5}
\]
şeklindedir.
\end{proof}

\begin{example}
Üç kenar uzunluğu verilen bir üçgenin geçerli olup olmadığını, geçerliyse eşkenar, ikizkenar veya çeşitkenar olduğunu belirleyen kodu yazınız.
\end{example}

\begin{proof}[Çözüm]
Kenar uzunlukları $a, b, c > 0$ olsun. Üçgen eşitsizliği gereği üçgenin çizilebilmesi için her kenar diğer iki kenarın toplamından kesinlikle küçük olmalıdır:
\[
a < b + c \quad \text{ve} \quad b < a + c \quad \text{ve} \quad c < a + b
\]
Eğer bu koşullar sağlanıyorsa üçgen geçerlidir ve türü belirlenebilir:
\begin{itemize}
    \item $a = b = c$ ise eşkenardır.
    \item Yalnızca iki kenar eşitse ikizkenardır: $(a = b) \lor (b = c) \lor (a = c)$.
    \item Aksi halde çeşitkenardır.
\end{itemize}

Program şu biçimde kurulabilir:
\begin{lstlisting}[language=CTemel]
#include <stdio.h>

int main(void) {
    long long a, b, c;
    if (scanf("%lld %lld %lld", &a, &b, &c) != 3) return 0;

    // Ucgen esitsizligi kontrolu
    if (a + b <= c || a + c <= b || b + c <= a) {
        printf("Gecersiz Ucgen\n");
    } else if (a == b && b == c) {
        printf("Eskenar\n");
    } else if (a == b || b == c || a == c) {
        printf("Ikizkenar\n");
    } else {
        printf("Cesitkenar\n");
    }

    return 0;
}
\end{lstlisting}
Kenar uzunluklarının büyük tam sayılar girilebileceği göz önüne alınarak \texttt{a + b} toplamında taşmayı önlemek amacıyla değişkenler \texttt{long long} tipinde tanımlanmıştır.
\end{proof}

\subsection{En Sık Yapılan Hatalar}
\begin{itemize}
    \item \textbf{Eşitlik operatörü \texttt{==} yerine atama operatörü \texttt{=} yazmak:}
    \texttt{if (x = 5)} yazıldığında $x$ değişkenine $5$ atanır. Atanan değer sıfırdan farklı olduğu için koşul her zaman doğru kabul edilir.
    \item \textbf{Matematiksel zincirleme karşılaştırmaları doğrudan C sözdizimine aktarmak:}
    Matematikteki $1 < x < 10$ ifadesi C dilinde \texttt{1 < x < 10} biçiminde yazılırsa sözdizimi hatası vermez ancak mantıksal olarak yanlış çalışır. C bu ifadeyi soldan sağa \texttt{(1 < x) < 10} şeklinde hesaplar. \texttt{1 < x} alt ifadesi duruma göre \texttt{0} veya \texttt{1} sonucunu üretir; bu değerler de her halükarda $10$'dan küçük olduğundan koşul $x$'ten bağımsız olarak daima doğru çıkar. Doğru yazım \texttt{1 < x \&\& x < 10} şeklindedir.
    \item \textbf{\texttt{switch} yapısında \texttt{break} komutunu unutmak:}
    İstenmeyen bir fall-through davranışı sonraki tüm durumların sırayla çalışmasına neden olur.
    \item \textbf{Koşul satırının sonuna noktalı virgül koymak:}
    \texttt{if (x > 0);} şeklinde bir noktalı virgül yerleştirildiğinde koşulun gövdesi boş sayılır; ardından gelen blok koşuldan bağımsız biçimde daima işletilir.
\end{itemize}
